% ===========================================================================
\section{Optimal input design}\label{sec:design}
% ===========================================================================

In distillation the labelling function is the teacher, and it can be queried at any prompt. The
training law $\cD$ is therefore a decision variable, not a given. \Cref{prop:covshift,prop:cov-term} bound
the covariate-shift term for a fixed $\cD$. This section asks which
$\cD$, possibly adaptive, minimises the deployment risk under $\mu$ at a fixed query budget, and what
the minimax rates are when outputs are sequences. \Cref{sec:design-answer} collects the answers and
\cref{sec:open} lists what remains open.

\subsection{Query model}

\begin{definition}[Designs and minimax deployment risk]\label{def:design}
A \emph{query} at a prompt $x$ returns a teacher sequence $y\sim\Pseq$ (\emph{hard labels}), or in
addition the soft labels $p_t(\cdot\mid s_t)$ along $y$ (\emph{soft labels}). A \emph{design}
$\cD\in\Delta(\cX)$ draws the prompts $x_1,\dots,x_n$ i.i.d.\ from $\cD$. An \emph{adaptive design}
$\mathfrak A$ chooses $x_i$ as a randomised function of the earlier queries and answers. After $n$
queries the learner outputs $\hat Q=(\hat Q^x)_{x\in\cX}$, any measurable function of the transcript.
We write $n_x$ for the number of queries at $x$. For a class $\mathcal P$ of teachers, the
\emph{minimax deployment risk} of $\cD$ is
\[
\mathfrak M_n(\mathcal P;\cD):=\inf_{\hat Q}\ \sup_{p\in\mathcal P}\ \E\,\cL_{\Hel^2}(\hat Q;\mu),
\qquad\cL_{\Hel^2}(\hat Q;\mu):=\E_{x\sim\mu}\Hel^2\big(\Pseq,\hat Q^x\big),
\]
where $\cL_{\Hel^2}$ extends \cref{def:risk} to arbitrary estimates. $\mathfrak M^{\mathrm{ad}}_n(\mathcal P)$
is the same quantity with a further infimum over adaptive designs.
\end{definition}
\noindent
We measure risk by squared Hellinger distance because it is bounded, symmetric and equivalent to
behavioural losslessness up to squaring: $\Hel^2\le\TV\le\sqrt2\,\Hel$ (\cref{prop:pinsker}). KL versions
are stated where they hold.

\subsection{What the covariate-shift terms can and cannot decide}\label{sec:design-shift}

Write $\rho(\mu\|\cD):=\sup_x\mu(x)/\cD(x)$ for the constant of \cref{as:shift}. Both constants in
\cref{prop:covshift}, $\rho(\mu\|\cD)$ and $\TV(\mu,\cD)$, are minimised exactly at $\cD=\mu$.
Indeed $\rho\ge1$, with equality iff $\cD=\mu$, because both laws have total mass one. Minimising the
shift terms therefore always returns $\cD=\mu$. But these are worst-case constants over \emph{all}
non-negative error profiles. The error profile of a trained student is constrained by its class, and
the constant that matters for design is the one restricted to that class.

\begin{proposition}[Class-restricted transfer and the design game]\label{prop:design-game}
\Pv\Ck Let $\cG$ be a set of functions $g:\cX\to[0,\infty)$ and define
\[
\kappa_\cG(\mu\|\cD):=\sup\Big\{\frac{\E_\mu g}{\E_\cD g}:\ g\in\cG,\ \E_\mu g>0\Big\},\qquad a/0:=\infty\text{ for }a>0 .
\]
\begin{enumerate}[label=(\roman*),nosep]
\item $\kappa_\cG(\mu\|\cD)\le\rho(\mu\|\cD)$, with equality when $\cG$ contains the indicators of all
singletons. \Cref{prop:covshift} is the case of the unrestricted class.
\item Let $\Theta_\TS$ be finite, $p=q_{\theta^\circ}$ with $\theta^\circ\in\Theta_\TS$, and
$\cG_\Theta:=\{x\mapsto\Hel^2(Q^x_\theta,Q^x_{\theta'}):\theta,\theta'\in\Theta_\TS\}$. The sequence-level
MLE of \cref{thm:mle}, fitted to $n$ hard-label queries i.i.d.\ from $\cD$, satisfies with probability
$\ge1-\delta$, for every horizon $T$,
\[
\E_{x\sim\mu}\Hel^2\big(\Pseq,Q^x_{\htheta}\big)\le\kappa_{\cG_\Theta}(\mu\|\cD)\,\frac{\log(|\Theta_\TS|/\delta)}n .
\]
\item Let $\cG^+:=\{g\in\cG:\E_\mu g>0\}$, on which $\kappa_\cG$ depends. If $\cX$ and $\cG^+$ are finite and
$\cG^+\ne\emptyset$, then $\kappa^\star_\cG:=\min_\cD\kappa_\cG(\mu\|\cD)=1/v$, where
\[
v:=\max_{\cD\in\Delta(\cX)}\ \min_{g\in\cG^+}\frac{\E_\cD g}{\E_\mu g}
=\min_{\lambda\in\Delta(\cG^+)}\ \max_{x\in\cX}\ \sum_{g\in\cG^+}\lambda_g\frac{g(x)}{\E_\mu g}\ \ge\ 1 .
\]
The optimal designs and the least-favourable mixtures $\lambda$ solve a pair of dual linear programs.
The design $\cD=\mu$ is optimal iff some $\lambda$ satisfies $\sum_g\lambda_g\,g/\E_\mu g\le1$ on all of $\cX$.
\item For every $\eps\in(0,1)$ there is an instance with two prompts and two students in which
$\kappa(\mu\|\mu)=1$, whereas $\min_\cD\kappa=\eps$ is attained at a point mass $\cD=\delta_{x_1}$, for which
$\rho(\mu\|\delta_{x_1})=\infty$.
\end{enumerate}
\end{proposition}
\begin{proof}
(i) $\E_\mu g\le\rho\,\E_\cD g$ for every $g\ge0$, and $g=\ind_{\{x\}}$ gives the ratio $\mu(x)/\cD(x)$.
(ii) By \cref{thm:mle}, $\E_\cD g\le\log(|\Theta_\TS|/\delta)/n$ for $g=\Hel^2(Q_{\theta^\circ},Q_{\htheta})\in\cG_\Theta$. If
$\E_\mu g>0$, then $\E_\mu g\le\kappa\,\E_\cD g$.
(iii) Put $\hat g:=g/\E_\mu g$. Then $1/\kappa_\cG(\mu\|\cD)=\min_g\E_\cD\hat g$, and maximising over $\cD$ is
a finite zero-sum game with payoff $\hat g(x)$. Von Neumann's minimax theorem gives the dual form, and LP
duality gives the programs. The design $\cD=\mu$ achieves $\min_g\E_\mu\hat g=1$, so $v\ge1$. Moreover $v=1$
iff the dual value is at most $1$, which is the stated condition.
(iv) Let $\cX=\{x_0,x_1\}$ and $\mu=(1-\eps,\eps)$, and take two students that agree at $x_0$ and differ at
$x_1$. The only non-zero profile in $\cG_\Theta$ is $g=\Hel^2(Q^{x_1}_{\theta_1},Q^{x_1}_{\theta_2})\ind_{\{x_1\}}$,
so $\kappa(\mu\|\cD)=\eps/\cD(x_1)$.
\end{proof}

\noindent
Part (iv) is the basic phenomenon. The best place to query can carry little or no deployment mass,
because that is where the candidate students disagree. The coefficient $\kappa_\cG$ plays the role of
the discrepancy distance and the transfer exponents of domain adaptation
\citep{mansour2009domain,hanneke2019value}. Under realisability the teacher is itself a student, so the
student class alone fixes $\cG_\Theta$, just as it fixes the discrepancy distance. What distillation adds
is that labels can be requested at any input, so $\cD$ is free. Bound (ii) inherits the union bound of \cref{thm:mle} and is not sharp for structured
classes. Sharp rates follow for the least structured class (\cref{sec:separable}) and, locally, for
smooth parametric classes (\cref{sec:local-design}).

\subsection{Unstructured teachers: minimax rates and the square-root law}\label{sec:separable}

\begin{definition}[Separable teachers]\label{def:separable}
$\mathcal P_{\mathrm{sep}}$ is the class of teachers with $\Pseq=p_x^{\otimes T}$, where $p_x\in\Delta(\cV)$ is
arbitrary and unrelated across prompts. Each prompt is an independent stochastic ``fact'', as in
\cref{cor:facts} but with stochastic answers. The student may be any law, so the problem is
realisable, and nothing learned at one prompt transfers to another.
\end{definition}

\begin{theorem}[Minimax rates for separable teachers]\label{thm:separable}
\Pv\Ck Let $k:=V-1$ be the per-prompt dimension, let $\mu_{(1)}\ge\mu_{(2)}\ge\cdots$ be the decreasing
rearrangement of $\mu$, and set
\[
S_n(\cD):=\sum_{x\in\cX}\mu(x)\min\Big\{1,\frac k{n\cD(x)}\Big\},\qquad
L_n:=\min_{m\ge0}\Big[\frac kn\Big(\sum_{i\le m}\sqrt{\mu_{(i)}}\Big)^2+\sum_{i>m}\mu_{(i)}\Big],
\]
with $k/0:=\infty$. Let $c_{\mathrm{sep}}:=3(1-e^{-1})/1024>1.8\cdot10^{-3}$. For every budget $n\ge1$ and every
horizon $T\ge1$:
\begin{enumerate}[label=(\roman*),nosep]
\item (any design) $c_{\mathrm{sep}}\,S_n(\cD)\le\mathfrak M_n(\mathcal P_{\mathrm{sep}};\cD)\le S_n(\cD)$;
\item (optimal design) $\inf_\cD S_n(\cD)=L_n$, attained by $\cD^\star\propto\sqrt\mu$ on the $m^\star$ most probable
prompts, where $m^\star$ attains the minimum in $L_n$ (if $n\le k$, then $S_n\equiv L_n=1$ and every design is
optimal);
\item (adaptivity) $c_{\mathrm{sep}}\,L_n\le\mathfrak M^{\mathrm{ad}}_n(\mathcal P_{\mathrm{sep}})\le L_n$, so adaptive
designs gain at most the constant factor $1/c_{\mathrm{sep}}$;
\item (the naive design) $S_n(\mu)=\sum_x\min\{\mu(x),k/n\}$.
\end{enumerate}
The upper bounds are achieved by the product Laplace estimator $\hat Q^x=\hat p_x^{\otimes T}$ with
$\hat p_x(a)=(N_x(a)+1)/(n_xT+V)$, where $N_x(a)$ counts token $a$ among the $n_xT$ tokens observed at $x$.
Deterministic allocations $(n_x)$ satisfy the bounds of (i) with $n\cD(x)$ replaced by $n_x$.
\end{theorem}
\begin{proof}[Proof sketch]
The full proof is in \cref{app:separable}.
\emph{Upper bound.} From $N$ draws, the Laplace estimator has $\E\KL(p\|\hat p)\le\log(1+\frac k{N+1})$, by
Jensen's inequality and $\E\frac1{1+\mathrm{Bin}(N,\pi)}\le\frac1{(N+1)\pi}$. With \cref{prop:pinsker}(iv) and
$\Hel^2(p^{\otimes T},q^{\otimes T})\le T\Hel^2(p,q)$, the risk at $x$ given $n_x$ is at most $\min\{1,k/(2n_x)\}$.
Averaging over $n_x\sim\mathrm{Bin}(n,\cD(x))$ gives $\min\{1,k/(n\cD(x))\}$.
\emph{Lower bound.} Assouad's lemma over sign perturbations $\pm\delta_x$ of $\lfloor V/2\rfloor$ token pairs at
each prompt, with $T\delta_x^2\le\frac14$. The triangle inequality for $\Hel$ projects any estimate onto the
cube. For adaptive designs the chain rule gives
$\KL(\P_\sigma\|\P_{\sigma'})=\E_\sigma[n_x]\,T\,\KL(p_{x,\sigma}\|p_{x,\sigma'})$ for vertices that differ only at prompt
$x$. So only the average allocation $\bar n_x$ enters, and $\sum_x\bar n_x\le n$. Because $\bar n_x$ depends on
$\delta$ under an adaptive design, $\delta$ is chosen as a fixed point of a continuous map (Brouwer).
\emph{Optimal design.} Rearranging $\cD$ to be ordered like $\mu$ cannot increase $S_n$, and after the
rearrangement the learned set $\{n\cD>k\}$ is a set of most probable prompts. Cauchy--Schwarz on
that set gives $L_n$, with equality for $\cD\propto\sqrt\mu$.
\end{proof}

\noindent
Three features of \cref{thm:separable} answer the design question for unstructured teachers. First,
the rate is \emph{uniform in the horizon}. A query returns $T$ tokens and the sequence-level loss is $T$
times more sensitive, and the two effects cancel exactly. Second, the optimal design is a
\emph{square-root law}: sample prompts with probability $\propto\sqrt{\mu(x)}$, which is flatter than $\mu$,
and learn only the head of $\mu$. Sampling from $\mu$ over-spends on frequent prompts and never reaches
the rare ones. Third, with hard labels adaptivity is worth at most a constant.

\begin{corollary}[Sampling from $\mu$ is rate-suboptimal for heavy tails]\label{cor:sep-rates}
\Pv\Ck As $n\to\infty$ with $k$ fixed:
\begin{enumerate}[label=(\alph*),nosep]
\item if $\mu$ is uniform on $K$ prompts, $S_n(\mu)\asymp L_n\asymp\min\{1,kK/n\}$, so $\cD=\mu$ is optimal up to
constants;
\item if $\mu_{(i)}\propto i^{-a}$ with $a>1$, then $S_n(\mu)\asymp(k/n)^{(a-1)/a}$, whereas $L_n\asymp(k/n)^{a-1}$ for
$a<2$, $L_n\asymp(k/n)\log^2(n/k)$ for $a=2$, and $L_n\asymp k/n$ for $a>2$;
\item if $\mu_{(i)}\propto r^i$ with $0<r<1$, then $S_n(\mu)\asymp(k/n)\log(n/k)$ and $L_n\asymp k/n$.
\end{enumerate}
In (b) the optimal design multiplies the exponent of the minimax rate by $\min\{a,a/(a-1)\}$, which
equals $a$ for $a\le2$.
\end{corollary}
\begin{proof}
(a) Direct. (b) $\mu_{(i)}\ge k/n$ iff $i\lesssim(n/k)^{1/a}$, and $\sum_{i>m}i^{-a}\asymp m^{1-a}$; this gives $S_n(\mu)$.
In $L_n$, $\sum_{i\le m}i^{-a/2}$ is $\asymp m^{1-a/2}$ for $a<2$ and $\asymp\log m$ for $a=2$. Balancing the two terms gives
$m^\star\asymp n/k$ and $m^\star\asymp n/(k\log(n/k))$ respectively. For $a>2$, $\sum_i\sqrt{\mu_{(i)}}<\infty$ gives
$L_n=O(k/n)$, and $L_n\ge\min\{1,k\mu_{(1)}/n\}$. (c) As in (b), with $\#\{i:\mu_{(i)}\ge k/n\}\asymp\log(n/k)$ and
$\sum_i\sqrt{\mu_{(i)}}<\infty$.
\end{proof}

\begin{corollary}[Soft labels: enumerate the head of $\mu$]\label{cor:sep-soft}
\Pv\Ck With soft labels, for every $n,T\ge1$,
\[
\tfrac14\sum_{i>n}\mu_{(i)}\ \le\ \mathfrak M^{\mathrm{ad}}_n(\mathcal P_{\mathrm{sep}})\ \le\ \sum_{i>n}\mu_{(i)},
\]
and the upper bound is achieved by querying the $n$ most probable prompts once each, a deterministic
non-adaptive allocation. For a design $\cD$, $\mathfrak M_n(\mathcal P_{\mathrm{sep}};\cD)$ lies between $\frac14$ and $1$
times $\sum_x\mu(x)(1-\cD(x))^n$. For $\cD=\mu$ this is the expected missing mass, which the Good--Turing
estimator estimates \citep{good1953population}. For $\mu_{(i)}\propto i^{-a}$ the optimal and naive rates are
$\asymp n^{1-a}$ and $\asymp n^{-(a-1)/a}$. For $\mu_{(i)}\propto r^i$ they are $\asymp r^n$ and $\asymp1/n$, and even the best
i.i.d.\ design only reaches $e^{-\Theta(\sqrt n)}$. With soft labels, enumeration therefore beats every i.i.d.\
design by a rate, while adaptivity adds nothing beyond enumeration.
\end{corollary}
\begin{proof}
One soft-label query at $x$ reveals $p_x$ (the soft label at $t=1$), hence $\Pseq$. For the upper bound,
output $\Pseq$ at queried prompts and any law elsewhere, where $\Hel^2\le1$. For the lower bound, draw each
$p_x$ independently and uniformly from $\{\delta_a,\delta_b\}$. Given the transcript, the teacher at an
unqueried prompt still has this prior. For every law $Q$ on $\cY$, the triangle inequality for $\Hel$ gives
$\frac12[\Hel^2(\delta_a^{\otimes T},Q)+\Hel^2(\delta_b^{\otimes T},Q)]\ge\frac14\Hel^2(\delta_a^{\otimes T},\delta_b^{\otimes T})=\frac14$. At
most $n$ prompts are queried, and under an i.i.d.\ design $x$ stays unqueried with probability $(1-\cD(x))^n$.
The rates follow as in \cref{cor:sep-rates}. In the geometric case the prompts with $\mu(x)\asymp1/n$
dominate the missing mass and contribute $\asymp1/n$. For the best i.i.d.\ design, take $m=\lceil\sqrt n\rceil$.
At least $m$ of the $2m$ most probable prompts have $\cD(x)\le1/m$, and each leaves mass
$\gtrsim r^{2m}(1-\frac1m)^n\ge r^{2m}e^{-2n/m}$ unqueried. Conversely, $\cD$ uniform on the top $m$ prompts leaves
at most $e^{-n/m}+r^m$.
\end{proof}

\begin{remark}[Unrestricted teachers, dimension, and KL]\label{rem:sep-general}
If $\Pseq$ may be any law on $\cY$, apply \cref{thm:separable} with vocabulary $\cY$ and horizon $1$. The same
rates hold with $k=V^T-1$, so the horizon now enters exponentially, through the per-prompt dimension.
Horizon-free rates therefore need structure, either within a prompt (here, position independence)
or across prompts (\cref{sec:local-design}). A bounded number of parameters is not enough. A smooth
one-parameter family $P_\theta=\bigotimes_t\mathrm{Bern}(\frac12+\gamma_t(\theta))$ whose curve $\gamma$ passes through all
$2^T$ points $\delta s$, $s\in\{\pm1\}^T$, contains that hypercube. Assouad's lemma over positions, with
$\delta^2\asymp1/n$, then gives minimax risk $\gtrsim\min\{1,T/n\}$ at a single prompt. What is needed is a
bound on local entropy that is uniform in $T$. For KL, deterministic allocations satisfy the upper
bound of \cref{thm:separable}(i) with the cost of an unlearned prompt raised from $1$ to at most $T\log V$.
Under an i.i.d.\ design a prompt of the head can also go unqueried, and the bound becomes
$\sum_x\mu(x)\big[\min\{T\log V,\frac{2k}{(n+1)\cD(x)}\}+T\log V\,(1-\cD(x))^n\big]$. The lower bound follows from
$\KL\ge2\Hel^2$. So the horizon enters the KL rate only through the mass left unqueried.
\end{remark}

\begin{heuristic}[The tempering law]\label{heur:tempering}
Suppose the risk at a prompt queried $m$ times decays like $m^{-\beta}$. Minimising
$\sum_x\mu(x)n_x^{-\beta}$ subject to $\sum_xn_x=n$ gives $n_x\propto\mu(x)^{1/(1+\beta)}$ on the prompts that are worth
learning:
\[
\cD^\star\propto\mu^{1/(1+\beta)}\quad\text{on the head of }\mu .
\]
The case $\beta=1$ ($\Hel^2$ or KL with hard labels) is the square-root law of \cref{thm:separable}. The limit
$\beta\to\infty$ (soft labels, where one query suffices) gives the uniform law on the head, as in
\cref{cor:sep-soft}. For behavioural TV, $\beta=\frac12$ because $\TV\asymp\Hel$ locally, and the law is $\mu^{2/3}$.
Temperature sampling of data sources, $q_i\propto p_i^\alpha$ with $\alpha$ between about $0.3$ and $0.7$, is a
standard heuristic in multilingual pre-training \citep{lample2019cross,xue2021mt5}. We quote this range from
memory; the primary sources were not reachable when we checked, so the exact values each paper uses
should be verified. The tempering law gives a minimax rationale for $\alpha=\frac12$ when each source is learned
at the parametric rate, and it predicts a smaller $\alpha$ for soft-label distillation.
\end{heuristic}

\subsection{Shared parameters: I-optimal design and the local minimax constant}\label{sec:local-design}

When the student's $k$ parameters are shared across prompts, a query at one prompt informs the student
everywhere. The design question then becomes a problem of classical optimal design
\citep{kiefer1960equivalence,pukelsheim1993optimal}, with a sequence-level information matrix. Let
$\Theta_\TS\subseteq\R^k$ be open and $p=q_{\theta^\circ}$, let $\cX$ be finite, and let $\theta\mapsto\log Q^x_\theta(y)$ be
twice continuously differentiable for all $x$ and $y$, with $Q^x_\theta(y)>0$ on a set that does not depend on
$\theta$. Softmax students qualify. The \emph{sequence Fisher information} of a prompt and of a design are
\[
\mathcal I_x:=\E_{y\sim\Pseq}\big[\nabla\log Q^x_\theta(y)\,\nabla\log Q^x_\theta(y)^\top\big]_{\theta=\theta^\circ}
=\sum_{t=1}^T\E_{s_t\sim\Pseq}\mathcal I(s_t),\qquad\mathcal I_\cD:=\E_{x\sim\cD}\mathcal I_x ,
\]
where $\mathcal I(s)=J(s)^\top(\diag p-pp^\top)J(s)$ is the per-state information of \cref{prop:rb}. The chain
rule holds because the per-token scores are martingale differences. Assume $\mathcal I_\mu\succ0$ and set
\[
\Phi(\cD):=\tr\big(\mathcal I_\cD^{-1}\mathcal I_\mu\big)\quad(\infty\text{ if }\mathcal I_\cD\text{ is singular}),\qquad
\Phi^\star:=\min_{\cD\in\Delta(\cX)}\Phi(\cD),\qquad\mathrm{lev}(x):=\tr\big(\mathcal I_\mu^{-1}\mathcal I_x\big).
\]
$\Phi$ is the I-optimality (integrated-variance) criterion for the sequence Fisher information, and
$\mathrm{lev}(x)$ is the Fisher leverage of prompt $x$, with $\E_\mu\mathrm{lev}=k$.

\begin{theorem}[Locally minimax input design]\label{thm:local-design}
\begin{enumerate}[label=(\alph*),leftmargin=*]
\item \Kn\Ck \emph{Fixed designs.} Let $\mathcal I_\cD\succ0$. The MLE $\htheta$ fitted to $n$ hard-label queries drawn
i.i.d.\ from $\cD$ satisfies
\[
n\,\cL_\KL(\htheta;\mu)\ \Rightarrow\ \tfrac12Z^\top\mathcal I_\mu^{1/2}\mathcal I_\cD^{-1}\mathcal I_\mu^{1/2}Z,\qquad Z\sim\cN(0,I_k),
\]
a limit law with mean $\frac12\Phi(\cD)$. At $\cD=\mu$ this is the $k/2n$ law of \cref{prop:exact}(b), now for
sequence-level KL: $\Phi(\mu)=k$.
\item \Pa{the regularity above} \emph{No adaptive design beats $\Phi^\star$.} For every adaptive design and
estimator,
\[
\lim_{b\to\infty}\ \liminf_{n\to\infty}\ \sup_{\|\theta-\theta^\circ\|_\infty\le b/\sqrt n}\ n\,\E_\theta\big\|\htheta_n-\theta\big\|^2_{\mathcal I_\mu}\ \ge\ \Phi^\star ,
\]
where $\E_\theta$ is taken with teacher $q_\theta$ and $\|v\|_A^2:=v^\top Av$. Since
$2\cL_\KL(\theta';\mu)=\|\theta'-\theta\|^2_{\mathcal I_\mu}(1+o(1))$ as $\theta'\to\theta$ when the teacher is $q_\theta$, this is the local
form of $n\,\cL_\KL\ge\frac12\Phi^\star$.
\item \Kn \emph{Adaptivity attains it.} Spend $\lceil\sqrt n\rceil$ queries on $\mu$ to obtain a pilot $\tilde\theta$, draw
the rest i.i.d.\ from a minimiser of $\Phi$ evaluated at $\tilde\theta$, and fit the MLE to all data. The limit
law of $n\,\cL_\KL(\htheta;\mu)$ then has mean $\frac12\Phi^\star$.
\item \Pv\Ck \emph{Equivalence theorem.} $\Phi$ is convex and the optimal information matrix is unique. A
design $\cD^\star$ is optimal iff
$\Psi(x):=\tr\big(\mathcal I_{\cD^\star}^{-1}\mathcal I_\mu\mathcal I_{\cD^\star}^{-1}\mathcal I_x\big)\le\Phi(\cD^\star)$ for all $x\in\cX$, and
then equality holds on the support of $\cD^\star$.
\item \Pv\Ck \emph{How much design can gain.} $\cD=\mu$ is optimal iff $\mathrm{lev}(x)\le k$ for every $x\in\cX$,
including prompts outside the support of $\mu$. In general
\[
\frac{k^2}{\max_x\mathrm{lev}(x)}\ \le\ \Phi^\star\ \le\ \big(\E_\mu\sqrt{\mathrm{lev}}\big)^2\ \le\ k ,
\]
and the importance-sampling design $\cD\propto\mu\sqrt{\mathrm{lev}}$ already achieves the middle bound. If the
parameters split into blocks, one of size $k_x$ for each prompt, then $\mathrm{lev}(x)=k_x/\mu(x)$,
$\Phi^\star=\big(\sum_x\sqrt{\mu(x)k_x}\big)^2$ and $\cD^\star\propto\sqrt{\mu(x)k_x}$. This is the square-root law of
\cref{thm:separable}.
\item \Pv\Ck \emph{Horizon.} $\Phi$ is unchanged when every $\mathcal I_x$ is multiplied by the same constant.
Since $\mathcal I_x=T\bar{\mathcal I}_x$, with $\bar{\mathcal I}_x$ the average per-position information along teacher paths
from $x$, the optimal design and the constant $\frac12\Phi^\star$ depend on $T$ only through the profile
$x\mapsto\bar{\mathcal I}_x$. The sequence-level rate is horizon-free; per token it is $\Phi^\star/(2nT)$.
\end{enumerate}
\end{theorem}
\begin{proof}
(a) The pairs $(x_i,y_i)$ are i.i.d.\ with Fisher information $\mathcal I_\cD$, because $\cD$ does not depend on
$\theta$. Hence $\sqrt n(\htheta-\theta^\circ)\Rightarrow\cN(0,\mathcal I_\cD^{-1})$ under standard regularity. The map
$\theta\mapsto\cL_\KL(\theta;\mu)$ has zero gradient and Hessian $\mathcal I_\mu$ at $\theta^\circ$; apply the delta method.
(b) The van Trees inequality applied to the whole transcript; see \cref{app:vantrees}.
(c) By (d) the optimal information matrix is unique, and it depends continuously on $\theta$. So the
information of the second stage converges to it. The MLE under a design whose empirical information
stabilises is asymptotically normal by the martingale central limit theorem
\citep{chaudhuri1993nonlinear}; \citet{pronzato2013design} give the general theory.
(d) $\cD\mapsto\mathcal I_\cD$ is affine, and $M\mapsto\tr(\mathcal I_\mu M^{-1})$ is strictly convex on positive definite
matrices. The directional derivative of $\Phi$ at $\cD$ towards $\delta_x$ is $\Phi(\cD)-\Psi(x)$, and a convex
function on the simplex is minimal iff all these derivatives are $\ge0$. Since $\E_{\cD^\star}\Psi=\Phi(\cD^\star)$,
equality must hold on the support \citep{kiefer1974general}.
(e) At $\cD=\mu$ we have $\Phi=k$ and $\Psi=\mathrm{lev}$; apply (d). For the upper bound, let $w=\mu/\cD$ and
$B:=\E_\mu[w\,\mathcal I_x]$. The block matrix $\E_\cD\big[(1,w)^\top(1,w)\otimes\mathcal I_x\big]$ has blocks $\mathcal I_\cD$,
$\mathcal I_\mu$, $\mathcal I_\mu$, $B$ and is positive semidefinite, so its Schur complement gives $\mathcal I_\cD\succeq\mathcal I_\mu B^{-1}\mathcal I_\mu$.
Hence $\Phi(\cD)\le\tr(\mathcal I_\mu^{-1}B)=\sum_x\mu(x)^2\mathrm{lev}(x)/\cD(x)$. Cauchy--Schwarz minimises the right side at
$\cD\propto\mu\sqrt{\mathrm{lev}}$, with value $(\E_\mu\sqrt{\mathrm{lev}})^2\le\E_\mu\mathrm{lev}=k$. For the lower bound, let
$M:=\mathcal I_\mu^{-1/2}\mathcal I_\cD\mathcal I_\mu^{-1/2}$. Then $\Phi(\cD)=\tr M^{-1}\ge k^2/\tr M$ and $\tr M=\E_\cD\mathrm{lev}\le\max_x\mathrm{lev}(x)$.
The block case is a direct computation.
(f) Immediate from the definition.
\end{proof}

\noindent
For smooth students, then, the answer is to query where the sequence Fisher information is large
relative to deployment mass, measured in the metric $\mathcal I_\mu$. Prompts whose teacher continuations
are long and uncertain carry more information, through the sum over $t$ in $\mathcal I_x$. The optimal design
can put mass outside the support of $\mu$, where \cref{prop:covshift} would treat the shift as infinite.

\begin{proposition}[No fixed design is uniformly optimal]\label{prop:adaptivity-gap}
\Pv\Ck For every $K\ge2$ and $\eps>0$ there is a one-parameter model with $K$ prompts in which every fixed
design satisfies $\sup_\theta\Phi_\theta(\cD)/\Phi^\star_\theta\ge K(1-\eps)$, the subscript marking the true parameter. By
\cref{thm:local-design}(c), the two-stage adaptive design attains the ratio $1$ at every $\theta$.
\end{proposition}
\begin{proof}
At prompt $j$ let every position emit $\mathrm{Bern}(\sigma(\theta c_j))$ independently, so that
$\mathcal I_j(\theta)=Tc_j^2\sigma'(\theta c_j)$. With $k=1$, $\Phi_\theta(\cD)/\Phi^\star_\theta=\max_j\mathcal I_j(\theta)/\E_\cD\mathcal I_j(\theta)$. Let $u^\star$
maximise $u^2\sigma'(u)$ and put $c_j=u^\star/\theta_j$ with $\theta_j=r^j$. At $\theta=\theta_j$,
$\mathcal I_i/\mathcal I_j=r^{2(j-i)}\sigma'(u^\star r^{j-i})/\sigma'(u^\star)$, which tends to $0$ as $r\to\infty$ for every $i\ne j$. So
the ratio at $\theta_j$ tends to $1/\cD(j)$, and $\max_j1/\cD(j)\ge K$.
\end{proof}

\begin{remark}[Token budgets, state queries and soft labels]\label{rem:state-design}
Because the teacher can be queried anywhere, a query can also be a single state $s=(x,y_{<t})$, which
returns one token or $p_t(\cdot\mid s)$. With hard labels, $n$ state queries from a law $\nu$ on states carry
information $n\,\E_\nu\mathcal I(s)$. Moreover $\mathcal I_\mu=T\,\E_{s\sim\bar d_\mu}\mathcal I(s)$, where
$\bar d_\mu:=\frac1T\sum_t\dteach t$ is taken with $x\sim\mu$. At equal token budgets, a sequence design $\cD$ is
therefore equivalent to the state design $\bar d_\cD$. State-level design optimises the same criterion over
the larger set of all laws on states, so it can only help. With soft labels in the realisable case the
Fisher information is infinite, and design reduces to identification. For a linear-softmax student,
soft labels at $d$ states with spanning features recover the teacher exactly (\cref{prop:exact}(a)). This
needs at least $\lceil d/T\rceil$ sequence queries, and exactly that many when each path contributes $T$ new
directions. Here the horizon helps.
\end{remark}

\subsection{Misspecified students: the design moves the target}\label{sec:mis-design}

\begin{proposition}[Unweighted designs inflate the approximation term]\label{prop:bias-design}
\Pv\Ck Let $e_\theta(x)\ge0$ be a per-input deployment divergence, and let
$\theta^\star_\cD\in\argmin_\theta\E_\cD e_\theta$ and $\theta^\star_\mu\in\argmin_\theta\E_\mu e_\theta$. Then
\[
\E_\mu e_{\theta^\star_\cD}\le\rho(\mu\|\cD)\,\rho(\cD\|\mu)\ \E_\mu e_{\theta^\star_\mu},
\]
and the factor is sharp: for any $\mu\ne\cD$ of full support on two prompts and any $\eta>0$, a class of
two students attains $(1-\eta)\,\rho(\mu\|\cD)\rho(\cD\|\mu)$.
\end{proposition}
\begin{proof}
By \cref{prop:covshift} twice and the optimality of $\theta^\star_\cD$,
\[
\E_\mu e_{\theta^\star_\cD}\le\rho(\mu\|\cD)\,\E_\cD e_{\theta^\star_\cD}\le\rho(\mu\|\cD)\,\E_\cD e_{\theta^\star_\mu}\le\rho(\mu\|\cD)\,\rho(\cD\|\mu)\,\E_\mu e_{\theta^\star_\mu}.
\]
For sharpness, label the prompts so that $\mu_1/\cD_1\ge1\ge\mu_2/\cD_2$. Then
$\rho(\mu\|\cD)\rho(\cD\|\mu)=\mu_1\cD_2/(\cD_1\mu_2)>1$. Take the profiles $e_{\theta_1}=(0,1)$ and $e_{\theta_2}=(s,0)$ with
$s=(1-\eta)\cD_2/\cD_1$. For small $\eta$, $\theta_2$ is optimal under $\cD$ and $\theta_1$ under $\mu$, and the ratio of their
$\mu$-risks is $s\mu_1/\mu_2$.
\end{proof}

\noindent
Under misspecification, then, the training law changes \emph{which} student is learned, not only how
well its risk transfers. Importance weighting removes this dependence. Minimising $\E_\cD[w\,\ell_\theta]$ with
$w=\mu/\cD$ targets $\theta^\star_\mu$ for every $\cD$ with $\mu\ll\cD$, and the design then affects only the variance.

\begin{proposition}[Optimal design for importance-weighted distillation]\label{prop:iw-design}
Let $\ell_\theta(x,o)$ be a per-query loss whose conditional mean given $x$ equals the deployment
divergence at $x$ up to a $\theta$-free term. Examples are $-\log Q^x_\theta(y)$ for hard labels and, by
\cref{thm:chain}, the soft-label cross-entropy summed along the teacher path. Let $\cX$ be finite, let
$R_\mu(\theta):=\E_\mu\E[\ell_\theta\mid x]$ be smooth with a unique interior minimiser $\theta^\star$, and let
$H:=\nabla^2R_\mu(\theta^\star)\succ0$. Put $S(x):=\E[\nabla\ell_{\theta^\star}\nabla\ell_{\theta^\star}^\top\mid x]$ and define the
\emph{sandwich leverage} $\mathrm{lev}_{\mathrm{sw}}(x):=\tr\big(H^{-1}S(x)\big)$.
\begin{enumerate}[label=(\alph*),nosep]
\item \Kn\Ck For $\mu\ll\cD$, the weighted M-estimator $\htheta_w$ satisfies
$n\big(R_\mu(\htheta_w)-R_\mu(\theta^\star)\big)\Rightarrow$ a limit law with mean $\frac12\E_\mu[(\mu/\cD)\,\mathrm{lev}_{\mathrm{sw}}]$
\citep{shimodaira2000improving}.
\item \Pv\Ck The optimal design is $\cD^\star\propto\mu\sqrt{\mathrm{lev}_{\mathrm{sw}}}$, with value
$\frac12(\E_\mu\sqrt{\mathrm{lev}_{\mathrm{sw}}})^2$, against $\frac12\E_\mu\mathrm{lev}_{\mathrm{sw}}$ for $\cD=\mu$. If $\mathrm{lev}_{\mathrm{sw}}\equiv c$ is
constant, then $\E_\mu[(\mu/\cD)\,\mathrm{lev}_{\mathrm{sw}}]=c\,(1+\chi^2(\mu\|\cD))$. With weighting, the covariate-shift
quantity that matters is therefore the leverage-weighted $\chi^2$ divergence, not $\rho$ or TV.
\item \Pv In the realisable case with hard labels, $S(x)=\mathcal I_x$, $H=\mathcal I_\mu$ and $\mathrm{lev}_{\mathrm{sw}}=\mathrm{lev}$,
so (b) is the importance-sampling design of \cref{thm:local-design}(e). The unweighted MLE does better
still, since $\Phi^\star\le(\E_\mu\sqrt{\mathrm{lev}})^2$: weighting only corrects a bias, and realisability
removes that bias. With soft labels, $S\equiv0$ in the realisable case, in line with \cref{rem:state-design}.
\end{enumerate}
\end{proposition}
\begin{proof}
(a) $\E_\cD[w\nabla\ell_{\theta^\star}]=\nabla R_\mu(\theta^\star)=0$, so $\sqrt n(\htheta_w-\theta^\star)\Rightarrow\cN(0,H^{-1}\Sigma H^{-1})$
with $\Sigma=\E_\cD[w^2S]$. A second-order expansion of $R_\mu$ at $\theta^\star$ gives the mean
$\frac12\tr(H^{-1}\Sigma)=\frac12\sum_x\mu(x)^2\mathrm{lev}_{\mathrm{sw}}(x)/\cD(x)$. Active learning with importance-weighted
estimators under misspecification is studied by \citet{kanamori2003active} and \citet{sugiyama2006active};
we have not re-verified whether they state this closed form.
(b) Cauchy--Schwarz, and $\sum_x\mu(x)^2/\cD(x)=1+\chi^2(\mu\|\cD)$.
(c) Substitute.
\end{proof}

\subsection{Where to query inside a sequence}\label{sec:path-design}

Because the teacher can be queried at any prefix, one could generate prefixes with a cheap proposal $r$,
for example the student, and ask the teacher only for soft labels. This moves the covariate shift from
prompts to states, and there it compounds with depth.

\begin{proposition}[State-level shift compounds with depth]\label{prop:path-shift}
\Pv\Ck Draw prompts from $\cD$ and prefixes from an autoregressive proposal $r$, so that the law of
depth-$t$ states is $d^{\,r}_t(x,y_{<t})=\cD(x)\prod_{j<t}r(y_j\mid s_j)$. The target law of forward-type risks
(\cref{thm:chain}) is $\dteach t$ with $x\sim\mu$. Let $w_t:=\dteach t/d^{\,r}_t$,
$\gamma:=\inf_s\max_ap(a\mid s)/r(a\mid s)\ge1$ and $\chi_0:=\inf_s\chi^2\big(p(\cdot\mid s)\,\|\,r(\cdot\mid s)\big)$.
\begin{enumerate}[label=(\roman*),nosep]
\item $\|w_t\|_\infty\ge\rho(\mu\|\cD)\,\gamma^{t-1}$ and $\E_{d^{\,r}_t}[w_t^2]\ge\big(1+\chi^2(\mu\|\cD)\big)(1+\chi_0)^{t-1}$. The constant
of \cref{prop:covshift} and the variance factor of importance weighting both grow exponentially with
depth.
\item For $r=p$ (teacher-generated prefixes), $w_t(x,y_{<t})=\mu(x)/\cD(x)$ for every $t$. Both quantities
keep their prompt-level values at every horizon.
\item The loss is not an artefact of weighting. Let $\cD=\mu$, and let $p(\cdot\mid s)\equiv p_0$ and $r(\cdot\mid s)\equiv r_0$
at depths below $t$, with $\mathsf K_0:=\KL(p_0\|r_0)>0$. Let the teacher's conditionals at depth-$t$ states
be unrelated to one another, as in \cref{def:separable}. A learner that receives soft labels at $N$
states drawn from $d^{\,r}_t$ has, in the worst case, expected depth-$t$ error
$\E_{s\sim\dteach t}\Hel^2(p(\cdot\mid s),\hat q(\cdot\mid s))\ge\frac14\big(1-\eps_t-Ne^{-(t-1)\mathsf K_0/2}\big)$ with $\eps_t\to0$. So $N$ must
grow like $e^{(t-1)\mathsf K_0/2}$.
\end{enumerate}
\end{proposition}
\begin{proof}
(i) Start at a prompt that maximises $\mu/\cD$ and follow the tokens that maximise $p/r$. For the second
moment, condition on $s_{t-1}$:
$\E_r[(p/r)^2(y_{t-1}\mid s_{t-1})\mid s_{t-1}]=1+\chi^2(p(\cdot\mid s_{t-1})\|r(\cdot\mid s_{t-1}))\ge1+\chi_0$. Induct down to
$\E_\cD[(\mu/\cD)^2]=1+\chi^2(\mu\|\cD)$.
(ii) The product over positions equals $1$.
(iii) Let $A_t$ be the set of depth-$t$ states with $\log\frac{P(y_{<t}\mid x)}{R(y_{<t}\mid x)}\ge(t-1)\mathsf K_0/2$. By the law
of large numbers $\dteach t(A_t)=1-\eps_t\to1$, and by a change of measure $d^{\,r}_t(A_t)\le e^{-(t-1)\mathsf K_0/2}$. The
expected number of queried states in $A_t$ is at most $Nd^{\,r}_t(A_t)$. So the expected $\dteach t$-mass of
unqueried states in $A_t$ is at least $1-\eps_t-Ne^{-(t-1)\mathsf K_0/2}$. At an unqueried state, the two-point prior
of \cref{cor:sep-soft} gives Bayes risk at least $\frac14$.
\end{proof}

\begin{remark}[Consequences for the design problem]\label{rem:path-design}
For forward-type risks (forward KL, and the teacher-prefix half of PA-JSD), teacher-generated prefixes
are the only state design whose shift constants do not grow with depth. The design freedom that
remains is the prompt law $\cD$, treated above. For reverse-type risks (\cref{eq:chain-rev}) the target is
the student's own prefix law $\dstud t$. The only horizon-free design is then on-policy, and it is
necessarily adaptive, because $\theta$ is being learned. On-policy distillation is adaptive input design,
and the DAgger-type guarantees of \cref{thm:pdl,rem:onpolicy} are guarantees for this design. Knowing the
labelling function makes $w_t$ computable, since $P(y_{<t}\mid x)$ follows from the teacher's soft labels
along the prefix, but it does not make $w_t$ small. This is the ``curse of horizon'' of off-policy
evaluation \citep{liu2018breaking}.
\end{remark}

\subsection{Answer to the design problem}\label{sec:design-answer}

\begin{center}\small
\begin{tabular}{@{}>{\raggedright\arraybackslash}p{0.19\linewidth}>{\raggedright\arraybackslash}p{0.21\linewidth}>{\raggedright\arraybackslash}p{0.22\linewidth}>{\raggedright\arraybackslash}p{0.3\linewidth}@{}}
\toprule
regime & optimal design & minimax deployment risk & adaptivity and horizon\\\midrule
shift bounds of \cref{prop:covshift} & $\cD=\mu$ & worst case over all error profiles & not a design criterion (\cref{prop:design-game})\\
finite realisable class & LP solution, possibly off the support of $\mu$ & $\le\kappa^\star_{\cG_\Theta}\log(|\Theta_\TS|/\delta)/n$ & horizon-free (\cref{prop:design-game})\\
separable, hard labels & $\propto\sqrt\mu$ on the head of $\mu$ & $\asymp L_n$ & adaptivity gains a constant; uniform in $T$ (\cref{thm:separable})\\
separable, soft labels & the $n$ most probable prompts, once & $\asymp\sum_{i>n}\mu_{(i)}$ & no gain from adaptivity (\cref{cor:sep-soft})\\
shared parameters, realisable & I-optimal for the sequence Fisher information & $\frac12\Phi^\star/n$, $\Phi^\star\le(\E_\mu\sqrt{\mathrm{lev}})^2\le k$ & adaptivity needed to reach $\Phi^\star(\theta^\circ)$, gains up to $|\cX|$, cannot beat it; horizon-free (\cref{thm:local-design,prop:adaptivity-gap})\\
misspecified, weighted & $\propto\mu\sqrt{\mathrm{lev}_{\mathrm{sw}}}$ & excess $\frac12(\E_\mu\sqrt{\mathrm{lev}_{\mathrm{sw}}})^2/n$ & unweighted designs move the target (\cref{prop:bias-design,prop:iw-design})\\
states inside sequences & teacher prefixes (forward), on-policy (reverse) & --- & mismatch costs $e^{\Theta(t\mathsf K_0)}$; on-policy is adaptive (\cref{prop:path-shift})\\\bottomrule
\end{tabular}
\end{center}

\noindent
In short: minimising the covariate-shift terms of \cref{prop:covshift} returns $\cD=\mu$, but the optimal
design samples the deployment law \emph{tempered by the square root of each prompt's information
leverage}, $\cD\propto\mu\sqrt{\mathrm{lev}}$, restricted to the head of $\mu$ that the budget can learn, with
teacher-generated continuations. For unstructured teachers $\mathrm{lev}=k/\mu$, so this design is $\sqrt\mu$ on
the head. It is then minimax optimal up to a constant (\cref{thm:separable}), and for heavy-tailed $\mu$
it improves the exponent of the rate (\cref{cor:sep-rates}). For shared parameters it is never worse
than $\mu$ and is within the factor $(\E_\mu\sqrt{\mathrm{lev}})^2/\Phi^\star$ of the I-optimal design
(\cref{thm:local-design}(e)). The minimax rates in the sequence setting are horizon-free whenever the
per-prompt teacher family has bounded dimension. Adaptivity changes constants, not rates, except
through the on-policy choice of prefixes. Experiment E4 (\cref{sec:experiments}) tests these
predictions, and \cref{sec:open} lists what remains open.

